\section*{Unit 7}
\addcontentsline{toc}{section}{Unit 7}
Tiga berkas berikut digunakan dalam Unit 7. Setiap berkas tetap mengikuti status hak atau lisensinya sendiri; tautan sumber dapat diklik dan tautan lisensi dicantumkan jika tersedia.
\begin{description}
\item[]\texttt{Stereographic} \texttt{projection} \texttt{in} \texttt{3D.png}
Mark.Howison (karya sendiri; domain publik).
\href{https://commons.wikimedia.org/wiki/File:Stereographic_projection_in_3D.png}{Halaman sumber di Wikimedia Commons}; status hak: Public domain.
\item[]\texttt{Manifold} \texttt{zahyou3.png}
132ninme di Wikipedia bahasa Jepang (karya sendiri; nama Unicode persis dicatat dalam manifes).
\href{https://commons.wikimedia.org/wiki/File:Manifold_zahyou3.png}{Halaman sumber di Wikimedia Commons}; lisensi \href{https://creativecommons.org/licenses/by-sa/3.0/}{CC BY-SA 3.0}.
\item[]\texttt{Circle} \texttt{-} \texttt{black} \texttt{simple.svg}
Dakdada (karya sendiri; domain publik).
\href{https://commons.wikimedia.org/wiki/File:Circle_-_black_simple.svg}{Halaman sumber di Wikimedia Commons}; status hak: Public domain.
\end{description}
Identitas revisi, ukuran, SHA-1 Wikimedia, SHA-256, URL berkas asli, dan hash turunan cetak tercatat dalam manifest media yang disertakan.
